\annexsubsection{\texorpdfstring{关于另一个二次型的二次型的迹 \(^{(1)}\)}{关于另一个二次型的二次型的迹 (1)}}

在本号中，记 E 为实向量空间，Q、H 为 E 上的两个正二次型。存在两个对称双线性型 \((x,y)\mapsto(x|y)_{\mathbf{Q}}\) 和 \((x,y)\mapsto(x|y)_{\mathbf{H}}\)，由下式刻画

\[
\mathrm{Q} (x) = (x | x) _ {\mathrm{Q}}, \quad \mathrm{H} (x) = (x | x) _ {\mathrm{H}}
\]

对于所有 \(x \in E\)。

称 Q 关于 H 的迹，并用 \(\operatorname{Tr}\left(\mathbf{Q}/\mathbf{H}\right)\) 表示如下定义的、可能有限也可能不有限的正实数：

\begin{enumerate}
\def\labelenumi{\alph{enumi})}
\item
  如果存在 \(x \in \mathbf{E}\)，使得 \(\mathrm{H}(x) = 0\) 且 \(\mathrm{Q}(x) \neq 0\)，则置 \(\operatorname{Tr}(\mathrm{Q} / \mathrm{H}) = +\infty\)。
\item
  在相反情况下，\(\operatorname{Tr}(\mathbf{Q}/\mathbf{H})\) 是形如 \(\sum_{i=1}^{p} \mathbf{Q}(e_i)\) 的数的集合的上确界，其中 \((e_1, \ldots, e_p)\) 遍历 E 中关于 H 正交规范的有限序列。
\end{enumerate}

设 E 是实希尔伯特空间，Q 是 E 上的正二次型。对所有 \(\mathrm{H}(x)=\|x\|^{2}\) 置 \(x\in E\)；于是 H 是 E 上的正二次型。如果 \(\mathrm{Tr}\left(\mathrm{Q}/\mathrm{H}\right)\) 有限，则称 Q 是核的。对于 E 中每个范数为 1 的 \(x\in E\)，有 \(\mathrm{Q}(x)\leqslant\mathrm{Tr}\left(\mathrm{Q}/\mathrm{H}\right)\)，从而 \(Q\leqslant\mathrm{Tr}\left(\mathrm{Q}/\mathrm{H}\right)\cdot\mathrm{H}\)；特别地，每个核型 Q 都连续。

备注. --- 1) 对于 E 的每个子空间 F，记 \(Q_{F}\) 为 Q 在 F 上的限制，记 \(H_{F}\) 为 H 在 F 上的限制。则 \(\mathrm{Tr}\left(\mathrm{Q}_{\mathrm{F}}/\mathrm{H}_{\mathrm{F}}\right) \leqslant \mathrm{Tr}\left(\mathrm{Q}/\mathrm{H}\right)\)，且 \(\mathrm{Tr}\left(\mathrm{Q}/\mathrm{H}\right)\) 是 \(\mathrm{Tr}\left(\mathrm{Q}_{\mathrm{F}}/\mathrm{H}_{\mathrm{F}}\right)\) 为有限维时各数 \(F \subset E\) 的上确界。

\begin{enumerate}
\def\labelenumi{\arabic{enumi})}
\setcounter{enumi}{1}
\tightlist
\item
  设 \(\mathbf{E}_1\) 是实向量空间，\(\mathbf{Q}_1\) 和 \(\mathbf{H}_1\) 是 \(\mathbf{E}_1\) 上的两个正二次型，\(\pi: \mathbf{E} \to \mathbf{E}_1\) 是满射线性映射。如果 \(\mathbf{Q} = \mathbf{Q}_1 \circ \pi\) 且 \(\mathbf{H} = \mathbf{H}_1 \circ \pi\)，则 \(\operatorname{Tr}(\mathbf{Q} / \mathbf{H}) = \operatorname{Tr}(\mathbf{Q}_1 / \mathbf{H}_1)\)。
\end{enumerate}

命题 1. --- 假设 E 是有限维的，且 H 非退化。

\begin{enumerate}
\def\labelenumi{\alph{enumi})}
\item
  存在  的自同态 \(u\)，由 （x, y 属于 \(\mathbf{E}\)）刻画。\((x|y)_{\mathbb{Q}} = (u(x)|y)_{\mathbb{H}}\)\(x, y\)\(\mathbf{E}\)
\item
  \(\operatorname{Tr}(\mathbf{Q} / \mathbf{H}) = \operatorname{Tr}(u)\)。
\item
  \(\operatorname{Tr}(\mathbf{Q} / \mathbf{H}) = \sum_{i=1}^{m} \mathbf{Q}(e_i)\)，其中 \((e_1, \ldots, e_m)\) 是 E 的任意关于 H 正交规范的基。
\end{enumerate}

 a) 由双线性型 \((x,y)\mapsto (x|y)_{\mathrm{H}}\) 非退化这一事实得到。E 中关于 H 正交规范的每个序列都可以补成 E 的关于 H 正交规范的一个基。因此，\(\operatorname {Tr}(\mathbf{Q} / \mathbf{H})\) 是 \(\sum_{i = 1}^{m}\mathbf{Q}(e_i)\) 形式的数在 E 的所有关于 H 正交规范的基 \((e_1,\ldots ,e_m)\) 上所成集合的上确界。要证明 b) 和 c)，只需证明对于每个这种基都有 \(\sum_{i = 1}^{m}\mathbf{Q}(e_i) =\) \(\operatorname {Tr}(u)\)。对于 \(a_{ij} = (u(e_i)|e_j)_{\mathrm{H}} = (e_i|e_j)_{\mathrm{Q}}\)，置 \(1\leqslant i,j\leqslant m\)；于是当 \(u(e_i) = \sum_{j = 1}^{m}a_{ij}e_j\) 时，\(1\leqslant i\leqslant m\)，从而

\[
\operatorname{Tr} (u) = \sum_ {i = 1} ^ {m} a _ {i i} = \sum_ {i = 1} ^ {m} \left(e _ {i} \mid e _ {i}\right) _ {\mathrm{Q}} = \sum_ {i = 1} ^ {m} Q \left(e _ {i}\right).\tag{Q.E.D.}
\]

命题 2. --- 假设 E 是有限维的。存在 E 的一个基 \((e_{1},\ldots,e_{n})\) 和一个满足 \(0\leqslant m\leqslant n\) 的整数 m，使得

\[
\mathrm{H} \left(\sum_ {i = 1} ^ {n} t _ {i} e _ {i}\right) = \sum_ {i = 1} ^ {m} t _ {i} ^ {2}\tag{1}
\]

对于实数 \(t_1, \ldots, t_n\)。此外，如果对于 \(\mathrm{H}(x) = 0\)，关系 \(\mathrm{Q}(x) = 0\) 蕴含 \(x \in \mathbf{E}\)，则 \(\operatorname{Tr}(\mathrm{Q}/\mathrm{H}) = \sum_{i=1}^{m} \mathrm{Q}(e_i)\)。

存在 E 的一个关于 H 正交的基 \((e_1', \ldots, e_n')\)。可以假定对基编号的方式使得当 \(\mathrm{H}(e_i') > 0\) 时 \(1 \leqslant i \leqslant m\)，而当 \(\mathrm{H}(e_i') = 0\) 时 \(m < i \leqslant n\)。于是，当 \(e_i = e_i' / \mathrm{H}(e_i')^{1/2}\) 时置 \(1 \leqslant i \leqslant m\)，当 \(e_i = e_i'\) 时置 \(m < i \leqslant n\)；关系 (1) 得到满足。

令 F 为由 \(e_{m+1}^{\prime},\ldots,e_{n}^{\prime}\) 生成的 E 的子空间；它正是满足 \(x\in E\) 的 \(\mathrm{H}(x)=0\) 的集合。记 \(\pi\) 为从 E 到 \(E_{1}=E/F\) 的典范映射。由于 Q 和 H 在 F 上为零，存在 \(Q_{1}\) 上的两个正二次型 \(H_{1}\) 和 \(E_{1}\)，使得 \(Q=Q_{1}\circ\pi\) 且 \(H=H_{1}\circ\pi\)。此外，\((\pi(e_{1}),\ldots,\pi(e_{m}))\) 是 \(E_{1}\) 的关于 \(H_{1}\) 正交规范的基，因而 \(H_{1}\) 非退化。

由命题 1 和备注 2，

\[
\operatorname{Tr} \left(\mathrm{Q} / \mathrm{H}\right) = \operatorname{Tr} \left(\mathrm{Q} _ {1} / \mathrm{H} _ {1}\right) = \sum_ {i = 1} ^ {m} \mathrm{Q} _ {1} \left(\pi (e _ {i})\right) = \sum_ {i = 1} ^ {m} \mathrm{Q} (e _ {i}).\tag{Q.E.D.}
\]

备注 3). --- 假设 E 有限维且 H 非退化。设 \((e_1, \ldots, e_n)\) 是 E 的一个基。置 \(q = ((e_i | e_j)_Q)_{1 \leqslant i, j \leqslant n}\)，\(h = ((e_i | e_j)_H)_{1 \leqslant i, j \leqslant n}\)。用命题 1 中的记号，u 关于 E 的基 \(u\)\((e_1, \ldots, e_n)\) 的矩阵等于 \(h^{-1}q\)，从而

\[
\operatorname{Tr} \left(\mathrm{Q} / \mathrm{H}\right) = \operatorname{Tr} \left(h ^ {- 1} q\right) = \operatorname{Tr} \left(q h ^ {- 1}\right).\tag{2}
\]

